\subsection{Relative safety and canonical extensions}
\label{app:core-relative-canonical}

A helper can be safe on another function's semantic budget while its own
canonical essential-source extension is not term-realizable.
Put
\[
 f=ABC\vee ACD\vee BCD\vee ABE\vee ACE\vee BCE,
 \qquad F=Z\land f.
\]
\begin{theorem}[Relative helpers need not be canonical]
\label{core:thm:relative-helper-not-canonical}
The extensions $D_f,D_F$ are term-realizable. A pure $S_F$-safe formula
separates $(0,AC)$ from $(0,BC)$, but no $R_F$-upset separating these
vertices has a term-realizable canonical extension. In particular,
canonical relative-safe helpers do not generate the bounded lattice of
relative-safe Boolean functions.
\end{theorem}
\begin{proof}
A canonical realizer of $f$ is
\[
 (A\vee B)\wedge\bigl(C(D\vee E\vee AB)\vee ABE\bigr).
\]
Writing $Y=AB$, the bracket is the paw $YE\vee C(D\vee Y\vee E)$.
Its canonicality follows from the cograph criterion and disjoint substitution
(\cref{core:thm:cograph-extension,core:prop:essential-disjoint-substitution}).
The two-threshold identity in \cref{core:prop:essential-threshold-band}
proves canonicality of the final meet. Disjoint conjunction with $Z$
makes $D_F$ realizable.

Set
\[
 H_0=AC\vee BCD\vee BCE\vee ABE,
 \qquad H_1=AC\vee CD\vee CE\vee ABE.
\]
The only $f$-zero vertex added by $H_0=f\vee AC$ is $AC$.
It has no outgoing initial deletion. A cube permitting deletion of $A$
would have lower endpoint contained in $C$ and upper containing $AC$ and
some $W\in\{B,D,E\}$; then $CW\to ACW$ is an $A$-pivotal edge inside it.
Deleting $C$ is symmetric. Every strict inclusion from $AC$ enters the
$f$-positive region, which cannot return to zero. Thus $H_0$ respects
$R_f^+$ and separates $AC,BC$.

We determine every $R_f^+$-upset $u$ separating them. It is nonzero,
so $u\geq f$, and monotonicity gives $u\geq H_0$. The following initial
rows of $f$ supply the indicated false-vertex routes:
\[
\begin{array}{c|c|c|c}
 L&U&j&f|_{[L,U]}\\\hline
 BD&ABCD&A&C\\
 \varnothing&ABDE&D&ABE\\
 D&ACDE&E&AC\\
 D&BCDE&E&BC
\end{array}
\qquad
\begin{aligned}
 ABD&\to BD\to B\subseteq BC,\\
 ADE&\to AD\subseteq ABD\to BD\to B\subseteq BC,\\
 BDE&\to BD\to B\subseteq BC.
\end{aligned}
\]
These and $BC$ are all maximal false vertices of $H_1$, so $u\leq H_1$.
The only undecided vertices are $CD,CE,CDE$. They are equivalent in $R_f$:
$CD\subseteq CDE\to CE$ uses the $D$-marked cube $[CE,ABCDE]$;
$CE\subseteq CDE\to CD$ uses the $E$-marked cube $[CD,ABCDE]$.
Both cofactors are $A\vee B$. Hence $u$ is exactly $H_0$ or $H_1$.

Neither canonical extension is realizable. Here are raw grouped-cut
certificates; every displayed cofactor is nonconstant and omits its mark:
\[
\begin{array}{c|c|c|c|c}
 &L&U&j&\text{cofactor}\\\hline
 H_0&\varnothing&ABC&B&AC\\
 &BD&ABCD&A&C\\
 &\varnothing&ABDE&D&ABE\\
 &CD&ABCDE&E&A\vee B\\
 &A&ABCDE&D&C\vee BE\\\hline
 H_1&B&ABDE&D&AE\\
 &D&ACDE&E&C\\
 &DE&ACDE&A&C\\
 &C&ABCDE&B&A\vee D\vee E
\end{array}
\]
For $H_0$ take the positive vertex $ABE$ and negative vertex $CDE$.
Its upper routes are $ABE\subseteq ABCDE\to ABCD\to ABC$ and
$ABE\subseteq ABDE$; its lower routes are
$A\subseteq ABD\to BD\to B\to\varnothing\subseteq CDE$ and $CD\subseteq CDE$.
The marks cover $ABE\setminus CDE=\{A,B\}$.
For $H_1$ take positive $ABE$ and negative $BC$.
Its upper routes are $ABE\subseteq ABDE$ and $ABE\subseteq ABCDE\to ACDE$;
its lower routes are $DE\to D\subseteq BD\to B\subseteq BC$ and $C\subseteq BC$.
The marks cover $ABE\setminus BC=\{A,E\}$.
Thus \cref{core:lem:grouped-cut} excludes both helpers. Distributing observer
union into pure components gives the same exclusion in the full signature:
every component would be safe and have the same Boolean restriction, since
disagreeing endpoint outputs would union to $\nv$.

By \cref{core:thm:conservation-disjoint},
$R_F=R_Z^+\times R_f^+$. A canonical $R_F$-upset separating the two
vertices would restrict at $Z=0$ to a canonical $R_f^+$-upset separating
$AC,BC$. It would be $H_0$ or $H_1$, a contradiction.

Finally an actual relative-safe separator exists. The read-once function
$g_0=C(A\vee B(D\vee E))$ respects $R_f$.
Indeed $C$ is essential on every nonconstant $f$-cube where $C$ is free: the
cofactors are $ABE$ and $AB\vee AD\vee BD\vee AE\vee BE$.
If $ABE$ has no true completion, a positive completion makes $C$ pivotal.
Otherwise, unless $A,B,E$ are all fixed one, set one free member to zero
and the other two to one; $C$ is again pivotal. The excluded case makes
$f$ constantly one. Thus $C$ never decreases in $R_f$.
Now $g_0=(f\wedge C)\vee AC$. At its $f$-positive winning vertices,
$f,C$ remain one along $R_f$; the only extra winning vertex is $AC$,
already treated above. Its canonical term is therefore $S_f$-safe.
Joining it with the canonical term for $f$ realizes $H_0$ safely on $S_f$.
Since $H_0\geq f$, this
formula is stably one on every $f$-constant-one cube, hence $S_F$-safe by
the cube classification for disjoint conjunction. It separates the stated
vertices despite its own nonrealizable canonical extension.
\end{proof}

\paragraph{Canonical intermediate values require unbounded gate arity.}
The difficulty also occurs inside a realizing circuit. A formula is
\emph{hereditarily canonical} if every subformula has a Boolean-valued
restriction and equals its canonical extension. We allow gates of arbitrary
finite arity, so this property records which intermediate functions are
retained after flattening associative gates. Even uniform-threshold gates
of bounded arity cannot provide such presentations for all realizable
canonical extensions.

A uniform-threshold gate denotes the canonical bichain extension of its
Boolean threshold. Fix $m\geq3$, write $[m]=\{1,\ldots,m\}$, use disjoint
variable sets $P=\{p_1,\ldots,p_m\}$ and $Q=\{q_1,\ldots,q_m\}$, and put
\[
 h_i=q_i\wedge\operatorname{Thr}_{m-2}(Q\setminus\{q_i\}),
 \qquad
 g_i=\bigvee_{j\ne i}p_j\vee h_i,
 \qquad F_m=\bigwedge_{i=1}^m g_i.
\]
Here a threshold on a variable set counts its true coordinates. Equivalently,
\[
 F_m=\operatorname{Thr}_2(P)\vee\bigwedge Q\vee\bigvee_i(p_i\wedge h_i).
\]
Indeed zero true $p$'s leave the conjunction of all $h_i$, which is
$\bigwedge Q$; exactly one true $p_i$ leaves $h_i$; two or more make
every factor true. Each $g_i$ has a pure canonical realizer $G_i$ by
uniform thresholds and disjoint substitution.

\begin{proposition}[An explicit relative-safe helper lattice]
\label{core:prop:canonical-helper-lattice}
The word $\bigwedge_iG_i$ realizes $D_{F_m}$. Write
$g_I=\bigwedge_{i\in I}g_i$ and $K=\bigvee_i g_i$. Its relative-safe
Boolean functions are exactly
\[
 \mathcal L_{S_{F_m}}
 =\{0,1,K\}\ \cup\ \{g_I:\varnothing\ne I\subseteq[m]\}.
\]
The interval $[F_m,K]$ is a Boolean lattice of rank $m$, with the
constant functions $0,1$ below and above it. In particular, $F_m$ has relative conservation.
Among these, exactly $0,1,K,F_m,g_1,\ldots,g_m$ have realizable canonical
extensions. Thus every conjunction of between two and $m-1$
distinct factors has a nonrealizable canonical extension.
For every realizer of $D_{F_m}$, its $m$ global Boolean cofactors fixing
one $p_i=1$ are relative-safe and generate it as a bounded lattice.
\end{proposition}
\begin{proof}
Write $V=P\cup Q$. Every nonconstant inessential marked cube is contained
in one of the following families. In the second row, $q_i\in R\subseteq Q$
and $|R|=m-2$; in the last two rows $j\ne i$.
\[
\begin{array}{c|c|c|c}
 \text{mark}&L&U&F_m|_{[L,U]}\\\hline
 p_i&Q\setminus\{q_i\}&V&\displaystyle\bigvee_{j\ne i}p_j\vee q_i\\
 q_i&\varnothing&P\cup R&\operatorname{Thr}_2(P)\\
 q_i&\{p_j\}&V\setminus\{q_j\}&\displaystyle\bigvee_{k\ne j}p_k\\
 q_i&\{p_j\}\cup(Q\setminus\{q_i,q_j\})&V&\displaystyle\bigvee_{k\ne j}p_k\vee q_j
\end{array}
\]
To check completeness, first mark $p_i$. No foreign $p$ is fixed one:
otherwise setting $p_i$ to one wins already at the lower assignment,
and inessentiality would make the original lower endpoint win.
If a foreign $p_j$ is available, set just that foreign $p$ to one.
The two $p_i$-cofactors are $h_j$ and $1$, so inessentiality forces
$h_j$ stably one. Thus at least $m-1$ $q$'s are fixed one.
They cannot all be fixed one on a nonconstant cube. If the unique missing
one were $q_k$ with $k\ne i$, setting it to zero and all foreign $p$'s
to zero would make $p_i$ pivotal. Hence all $Q\setminus\{q_i\}$ are
fixed one. If no foreign $p$ is available, the winning upper endpoint
and inessentiality require all $q$'s available; the same pivotal
assignments again force every foreign $q$ fixed one. This is the first family.

Now mark $q_i$. At most one $p$ is fixed one. Its pivotal conditions,
according to the true $p$-coordinates, are
\[
\begin{array}{c|c}
 \text{true }p\text{'s}&\text{$q_i$ pivotal}\\\hline
 \varnothing&\displaystyle\bigwedge_{j\ne i}q_j\\
 \{p_i\}&\operatorname{Thr}_{m-2}(Q\setminus\{q_i\})\\
 \{p_j\},\ j\ne i&q_j\ \wedge\ [\text{exactly }m-3\text{ true in }Q\setminus\{q_i,q_j\}]\\
 \text{at least two}&0
\end{array}
\]
If no $p$ is fixed one, the first condition excludes an upper endpoint
containing all of $Q$. If that endpoint contains $m-1$ $q$'s, let $q_k$
be the missing one. The other two conditions exclude every available
$p_j$ except possibly $p_k$, which alone cannot make the upper endpoint
win. Thus at most $m-2$ $q$'s are available, giving the second family.
If $p_i$ is fixed one, the second condition gives the same conclusion.
If instead $p_j$ is fixed one, put $R'=Q\setminus\{q_i,q_j\}$.
Avoiding the third condition means that $q_j$ is unavailable, all of
$R'$ is fixed one, or at most $m-4$ elements of $R'$ are available.
These give the third, fourth, and second families, respectively.

On the first family $G_i=\nv$ and all other $G_j=c$.
On the second every $G_i=\nv$. On either remaining family $G_j=\nv$
and the other helpers equal $c$. Restricting an unmarked neutral input
to an endpoint cannot introduce $\pv$, so all $G_i$ are $S_{F_m}$-safe.
Their conjunction has Boolean function $F_m$ and is safe on its entire
semantic budget. Stable corners and essential marked sources prove that
it is $D_{F_m}$.

The initial preorder $R_{F_m}$ has a bottom class
\[
 E_0=\{S\subseteq Q:|S|\leq m-2\},
\]
$m$ incomparable middle classes
\[
 E_i=\{\{p_i\}\cup S:S\subseteq Q,\ |S|\leq m-2\}
       \ \cup\ \{Q\setminus\{q_i\},\ \{p_i\}\cup(Q\setminus\{q_i\})\},
\]
and a top class consisting of every $F_m$-positive vertex.
The displayed Boolean formula classifies all negative vertices into
these classes. The second row family clears small $Q$-sets, with or
without $p_i$. The third clears $Q\setminus\{q_i\}$ in the presence
of $p_i$, and the first identifies this $Q$-set with $p_i$.
These give the stated negative classes.

Every positive vertex reaches $V$ by inclusion. Conversely, $V$ reaches
every minterm. To reach $p_ip_j$, delete the other $p$'s by the first
family, then $q_i,q_j$ by the fourth, and all remaining $q$'s by the
second. To reach $Q$, delete all $p$'s by the first. To reach
$\{p_i\}\cup(Q\setminus\{q_k\})$ for $k\ne i$, delete the other
$p$'s by the first, then $q_k$ by the fourth. These are all the minterms.
Inclusion puts $E_0$ below each $E_i$ and each $E_i$ below the positive
class. The safe features $g_i$ exclude further comparisons: their flag
vectors are $0^m$ on $E_0$, one everywhere except coordinate $i$ on
$E_i$, and $1^m$ on the positive class.

The upsets of this preorder are exactly the functions in the statement.
All of them have safe formulas, being monotone combinations of the
$G_i$ and constants. Hence they are exactly $\mathcal L_{S_{F_m}}$.
Moreover $F_m|_{p_i=1}=g_i$. Since canonical extension commutes with
Boolean cofactors, applying this assignment to any realizer of $D_{F_m}$
gives $D_{g_i}$. These cofactors therefore supply the same safe generators
independently of the original word.
The function $K=\bigvee P\vee\operatorname{Thr}_{m-1}(Q)$ has a canonical
realizer by disjoint substitution.

Suppose $2\leq|I|\leq m-1$, and choose $i,j\in I$ and $k\notin I$.
Keep the three paired indices $i,j,k$, fix all other $p$'s to zero and
all other $q$'s to one. The retained factors $g_i,g_j$ become their
six-variable versions at $m=3$. Every other factor in $I$ becomes
$\bigvee_{\ell\in\{i,j,k\}}p_\ell\vee
\operatorname{Thr}_2(\{q_i,q_j,q_k\})$, which dominates both retained
factors. Finally fix $p_k=0$. The resulting cofactor is
\[
 p_ip_j\vee q_iq_j\vee q_k(p_iq_i\vee p_jq_j),
\]
the noncanonical square switch of
\cref{core:thm:canonical-cofactor-obstruction}, after renaming.
Canonical extensions commute with Boolean cofactors, proving the claimed
nonrealizability of every intermediate $g_I$.
\end{proof}

\begin{corollary}[Unbounded arity of canonical gate presentations]
\label{core:cor:canonical-gate-arity}
For every $m\geq3$, the minimum possible maximum gate arity of a
hereditarily canonical formula for $D_{F_m}$ is exactly $m$, even when
arbitrary uniform-threshold gates and observer union are allowed.
The upper bound is attainable using only AND and OR gates. The same
minimum holds for circuits.
\end{corollary}
\begin{proof}
Interpret a uniform-threshold gate by its canonical bichain extension.
At a gate with Boolean restriction $F_m$, each $S_{F_m}$ row makes the
outer threshold nonconstant on the independent cube of its children's
endpoint values: its lower corner gives $F_m(L)=0$ and its upper corner
gives $F_m(U)=1$. Every free argument of a nonconstant uniform threshold
is essential. Thus a marked child would give a marked parent, and all
children are $S_{F_m}$-safe.

By the proposition, a canonical child has Boolean function $0$, $1$,
$K$, $F_m$, or one of the $g_i$. If a child has restriction $F_m$,
descend to it. Otherwise each child agrees on the positive class
and on every middle class except possibly its one index $i$ when the
child is $g_i$. Fewer than $m$ children therefore leave some $E_i$
with exactly the same child-value vector as the positive class. No gate
can then output $F_m$. Finite descent cannot finish at a variable or
constant, so some gate has arity at least $m$.
In a hereditarily canonical formula an observer-union gate has identical
Boolean restrictions in all children; their canonical extensions agree,
so it can be erased without increasing the other gate arities.

For the upper bound use the flat $m$-ary conjunction of the $G_i$.
Within $G_i$, the threshold on $m-1$ variables is a flat disjunction of
its $m-1$ monomials, each an $(m-2)$-ary conjunction. Every one of these
subformulas is canonical, as are the disjoint conjunction with $q_i$
and the disjoint disjunction with the other $p$'s. All remaining gate
arities are at most $m-1$. Finally, unfolding a circuit into a formula
preserves both gate arities and hereditary canonicality, so sharing
cannot improve the minimum.
\end{proof}
